\chapter{Kritische Bewertung und Praxisrelevanz}

\section{Governance-Aspekte: Einbindung in die \ac{IT}-Infrastruktur}
Wichtig für eine erfolgreiche Einbindung eines Maskierungsalgorithmus ist die Vermeidung von sogenannten Schatten-\acp{KI}. Darunter versteht man \ac{LLM}-Systeme, welche nicht von einer \ac{IT}-Abteilung bereitgestellt wurden und somit auch keinen Maskierungsalgorithmus verwenden \footcite{mansner2025shadowai}.
Da Mitarbeiter bei fehlenden oder unperformanten Systemen oft stattdessen auf nicht geschützte \ac{LLM}-Systeme zurückgreifen, ist es wichtig, interne \ac{LLM}-Systeme bereitzustellen, die nicht nur aktiv beworben werden, sondern auch auf dem aktuellen Stand der Technik sind.

Der hybride Ansatz ist hierbei entscheidend, da er die hohe Präzision von \ac{NER} mit der minimalen Latenz von \ac{Regex} kombiniert. Nur eine performante Lösung, die den Arbeitsfluss nicht durch lange Wartezeiten unterbricht, kann die Akzeptanz der Mitarbeiter sicherstellen und somit den Rückzug in ungeschützte Schatten-KI-Systeme effektiv verhindern \footcite{naseer2025recap}.

Zudem sollte man abwägen, ob ein lokales \ac{LLM}-Modell (on-premise), zum Beispiel auf eigenen Servern, günstiger und effizienter wäre, als ein pseudonymisierendes externes \ac{LLM}-Modell, welches über eine \ac{API} angesprochen wird \footcite{costbenefit2025onprem}. Dieses Modell würde keinen Maskierungsalgorithmus verwenden, was die Antwortqualität und Latenz verbessern würde.

Zuletzt ist es auch wichtig, dass Mitarbeiter über „Residual Risks“ (Restrisiken) geschult werden, denn kein Maskierungsalgorithmus kann die absolute Sicherheit der Daten garantieren \footcite{staab2025sokprivacyparadox}. Deshalb ist es wichtig, nicht nur die technische Implementierung, sondern auch die organisatorischen Maßnahmen zu prüfen.

[Image of cloud vs on-premise LLM infrastructure comparison]

\section{Compliance: \ac{DSGVO} \& Privacy by Design}
Auch die \ac{DSGVO} nennt explizit die datenschutzfreundliche Voreinstellung und den Datenschutz durch Technik (Privacy by Design) als Grundsätze der Datenverarbeitung. In Artikel 25 der \ac{DSGVO} heißt es dazu:
\begin{quote}
  Der Verantwortliche trifft geeignete technische und organisatorische Maßnahmen, um sicherzustellen, dass grundsätzlich nur Datenverarbeitungen vorgenommen werden, die für den jeweiligen Zweck erforderlich sind.
\end{quote}
Gleichzeitig wird in Artikel 32 der \ac{DSGVO} sogar explizit die Pseudonymisierung als Maßnahme genannt, welche zur Erfüllung dieser Pflicht beitragen können.
Zusätzlich wird in dem neuen EU AI Act die Verwendung von \ac{KI}-Systemen zur Analyse von biometrischen Daten oder zur Erkennung von Emotionen stark eingeschränkt, was zeigt, dass die Bedeutung von Datenschutz und Datensicherheit in der EU zunimmt \footcite{europeanparliament2025interplayaact}.